\documentclass{kothesis} % クラスファイルにまとめました
\usepackage{preamble} % 適宜プリアンブルをpreamble.styに追加
% 表紙
\affiliation{東京大学 工学部 電気電子工学科}
\thesistype{2018}{解説資料}
\title{クラスファイル ``kothesis'' の使い方 \\(まだ未完成です。.texファイルをみて察してください)}
% 指導教員は2人まで指定可能
\professor{指導教員名}{教授}
\date{yyyy年mm月dd日提出}
\author{00-000000}{電電 太郎}

\begin{document}
    %タイトルページを作成
    \maketitle
    % 概要を追加
    \includeabstract{chapter/abstract}
    % 目次，図目次，表目次を追加
    \makemokuji

    % チャプタごとに別ファイルに
    \include{chapter/introduction}
    \include{chapter/setup}
    \include{chapter/modeling}

%    \part{理論} % 2部構成の場合

    \include{chapter/conventional}
    \include{chapter/proposed}

%    \part{応用}

    \include{chapter/simulation}
    \include{chapter/experiment}
    \include{chapter/conclusion}

    % 謝辞，付録，発表文献は特殊なチャプタです
    \includeacknowledgments{chapter/acknowledgments}
    \includeappendix{chapter/appendix}

    % 参考文献をBibTeXで
    \bibliographystyle{IEEEtran}
    \bibliography{ref}

    \includepublication{chapter/publication}

    \printindex

\end{document}
